\chapter{Drafting the Model}\label{chap:3}

\begin{quote}
\emph{``All models are wrong, but some are useful.''} --- George E. P.
Box
\end{quote}

There are two fundamental challenges to overcome before implementing an
optimization model in a solver. The first is to express the problem in a
precise form. The second is to translate that form into something the
interface of your chosen solver can understand. Further challenges will
follow later, especially making the model run fast, but this chapter
focuses only on the first two.

The previous chapter gave us the notation needed to start sketching a
model. In this chapter, we look at a structured way to do so. The
structure is not linear: for a complicated problem, you usually identify
the core, formulate it, get it into a runnable form, and then extend it
with the remaining rules and objectives.

The five modeling ingredients in this chapter --- entities, parameters,
variables, constraints, and objective --- are therefore not a checklist
that you complete once. They are a recurring guide. New constraints may
require new variables; a new objective term may change which parameters
matter; a newly discovered entity may force you to restructure earlier
parts of the formulation. With experience, the iterations tend to become
fewer and smaller.

\section{The Five Components}\label{sec:3.1}

A complete optimization model has five components that provide structure
for thinking and writing:

\begin{enumerate}
\item
  \textbf{Entities} --- what objects does the problem talk about? They
  become \emph{sets}.
\item
  \textbf{Parameters} --- what data is known up front? They become
  \emph{constants indexed by the sets}.
\item
  \textbf{Variables} --- what is the solver allowed to choose? They
  become \emph{unknowns with a domain}.
\item
  \textbf{Constraints} --- what makes a choice valid? They become
  \emph{relations on variables}.
\item
  \textbf{Objective} --- what makes one valid choice better than
  another? It becomes \emph{one expression with sense
  \lstinline!min! or \lstinline!max!};
  multi-objective optimization is a topic of its own.
\end{enumerate}

The order is a guideline. Entities and parameters define the language;
variables depend on that language; constraints relate the variables; the
objective evaluates valid choices.

For the Knapsack problem from the previous chapter, the five components
read:

\begin{enumerate}
\item
  \textbf{Entities.} The items, \(I\).
\item
  \textbf{Parameters.} A value \(c_i \in \mathbb{R}_{\ge 0}\) and a
  weight \(w_i \in \mathbb{R}_{\ge 0}\) for every \(i \in I\), plus a
  knapsack capacity \(C \in \mathbb{R}_{\ge 0}\).
\item
  \textbf{Variables.} One binary variable per item, \(x_i \in \{0, 1\}\)
  for \(i \in I\), with \(x_i = 1\) meaning \emph{take item \(i\)}.
\item
  \textbf{Constraints.} The total weight must not exceed capacity:
  \(\sum_{i \in I} w_i x_i \le C\).
\item
  \textbf{Objective.} Maximize total value:
  \(\max \sum_{i \in I} c_i x_i\).
\end{enumerate}

Now suppose we want to integrate the scenario-based worst-case objective
from the previous chapter. We have to change a few components, but the
model still looks similar:

\begin{enumerate}
\item
  \textbf{Entities.} Items \(I\) and scenarios \(S\).
\item
  \textbf{Parameters.} A per-scenario value
  \(c_i^s \in \mathbb{R}_{\ge 0}\) for every \(i \in I\) and
  \(s \in S\); weights \(w_i\) and capacity \(C\) as before.
\item
  \textbf{Variables.} Unchanged: \(x_i \in \{0, 1\}\) for \(i \in I\).
\item
  \textbf{Constraints.} Unchanged: \(\sum_{i \in I} w_i x_i \le C\).
\item
  \textbf{Objective.} Maximize the worst-case value:
  \(\max \min_{s \in S} \sum_{i \in I} c_i^s x_i\).
\end{enumerate}

This is a precise paper formulation, but it is not yet in a form
suitable for most solvers. The objective contains a \(\min\) over
scenarios inside the maximization, and most solvers will not accept that
nested expression as written. The standard remedy is an
auxiliary-variable reformulation: introduce a variable \(z\) for the
worst-case value and bound it from above by every scenario's value.

The solver-facing version is:

\begin{enumerate}
\item
  \textbf{Entities.} Items \(I\) and scenarios \(S\).
\item
  \textbf{Parameters.} \(c_i^s\), \(w_i\), and \(C\).
\item
  \textbf{Variables.} \(x_i \in \{0, 1\}\) for \(i \in I\), and an
  auxiliary variable \(z\) representing the worst-case value.
\item
  \textbf{Constraints.} The capacity constraint \[
  \sum_{i \in I} w_i x_i \le C,
  \] and the scenario-value constraints \[
  z \le \sum_{i \in I} c_i^s x_i \qquad \forall s \in S.
  \]
\item
  \textbf{Objective.} Maximize \(z\).
\end{enumerate}

Maximizing \(z\) subject to one inequality per scenario forces \(z\) up
to the smallest scenario value, recovering the original
\(\min_{s \in S}\) exactly.

\section{A Worked Example: Hospital Ward Rostering}\label{sec:3.2}

Consider the following problem. A small hospital ward needs to schedule
its nursing staff for the coming week. The ward has a fixed list of
shifts to fill. Each shift has a date, a required skill --- say,
\emph{ICU}, \emph{pediatrics}, or \emph{general care} --- and a
headcount specifying how many nurses must be on duty. The ward employs a
roster of nurses, each of whom holds a personal set of skills and is
available on certain days but not others. Assigning a particular nurse
to a particular shift incurs a cost that captures wage rates, shift
differentials, and union overtime rules.

The ward manager wants to staff every shift, respect availability and
skills, and spend as little as possible.

We will now extract the five components. Unlike in the knapsack example,
they are no longer almost forced by the problem statement. We have
several kinds of objects, several indexed data tables, and constraints
that refer to different combinations of workers, shifts, days, and
skills. The point of the five-component structure is to keep that
complexity visible and organized.

\subsection{Step 1 --- Entities become sets}\label{step-1-entities-become-sets}

Read the problem statement and underline every noun that names a class
of objects. \emph{Nurses, shifts, days, skills.} Each class becomes a
set.

\begin{symboltable}{lL}
Symbol & Description \\ \midrule
\(W\) & Workers (nurses) --- the staff to be scheduled. \\
\(S\) & Shifts --- the work blocks to be filled this week. \\
\(T\) & Days --- the calendar dimension, \(T = \{1, \dots, 7\}\). \\
\(K\) & Skills --- the qualifications the ward tracks. \\
\end{symboltable}


A set is the \emph{type} of an index. Once we have written \(w \in W\)
and \(s \in S\), we can index any quantity in the model by \((w, s)\)
--- variables, parameters, constraints, neighborhoods.

\subsection{Step 2 --- Parameters are the input data}\label{step-2-parameters-are-the-input-data}

Anything known before we solve is a parameter. Parameters are often
indexed by the sets from step 1, and each one should carry a domain ---
the values it may take --- together with an informal description that
fixes what the symbol means. Some parameters are raw input data; others
are derived from that input to make the model easier to write.

\begin{symboltable}{lllL}
Symbol & Indices & Domain & Description \\ \midrule
\(c_{w,s}\) & \(w \in W,\; s \in S\) & \(\mathbb{R}_{\ge 0}\) & Cost of assigning worker \(w\) to shift \(s\). \\
\(a_{w,s}\) & \(w \in W,\; s \in S\) & \(\{0, 1\}\) & Availability flag: \lstinline!1! iff \(w\) is available for \(s\). \\
\(r_s\) & \(s \in S\) & \(K\) & Skill required by shift \(s\). \\
\(h_w\) & \(w \in W\) & \(2^K\) & Set of skills held by worker \(w\). \\
\(n_s\) & \(s \in S\) & \(\mathbb{N}_0\) & Number of workers needed on shift \(s\). \\
\(S_t\) & \(t \in T\) & \(2^S\) & Subset of shifts that fall on day \(t\) (derived from shift dates). \\
\end{symboltable}


The domains are explicit. \(c_{w,s}\) is a non-negative real,
\(a_{w,s}\) is binary, \(n_s\) is a non-negative integer, and \(h_w\) is
a subset of \(K\) (so its domain is the power set \(2^K\)). Each shift
requires one skill, while each worker may hold multiple skills --- a
detail that matters once we write the skill-match constraint.

The last row illustrates that a parameter need not be raw input. \(S_t\)
is derived from the dates of the shifts: for each day \(t\), it collects
exactly those shifts that occur on that day. Including such derived
parameters in the table is useful because it gives later constraints a
clean vocabulary. Instead of repeatedly writing ``all shifts that fall
on day \(t\),'' we can simply write \(S_t\).

\subsection{Step 3 --- Variables are the unknowns}\label{step-3-variables-are-the-unknowns}

This is the step where the problem stops being only a description and
becomes a model. Each decision unit becomes one variable, declared with
three things: a name, an index range, and a domain. What counts as a
single decision unit is itself a modeling choice: the same problem can
often be modeled at different granularities, with real consequences for
solver performance.

For the rostering problem, the natural first decision unit is \emph{does
worker \(w\) take shift \(s\)?} --- one binary choice per worker--shift
pair.

\begin{symboltable}{lllL}
Symbol & Indices & Domain & Description \\ \midrule
\(x_{w,s}\) & \(w \in W,\; s \in S\) & \(\{0, 1\}\) & \lstinline!1! iff worker \(w\) is assigned to shift \(s\), else \lstinline!0!. \\
\end{symboltable}


That single line declares \(|W| \cdot |S|\) variables --- one per cell
of the worker--shift matrix. Many of those combinations may be trivially
infeasible: the worker may be unavailable, or may lack the skill
required by the shift. In the simple formulation, we still declare the
full matrix and let constraints rule out the invalid combinations.

One could prune the index set up front and create variables only for
worker--shift pairs that are possible in principle. That is often a
useful implementation choice, but it is not needed for the first paper
model. At this stage, the full matrix is easier to read: every potential
assignment has a variable, and every rule appears explicitly as a
constraint. If an implementation detail is important enough to remember,
record it as a note rather than complicating the first formulation.

The decision here is \emph{take or leave}, so the natural variable is
binary. Other problems suggest other shapes:

\begin{symboltable}{lL}
Decision phrasing & Variable shape \\ \midrule
``Take it or not?'' & \(x_i \in \{0, 1\}\) --- binary \\
``How many?'' & \(x_i \in \mathbb{N}_0\) --- non-negative integer \\
``When does it start?'' & \(\text{st}_i \in [0, H]\) --- time variable, bounded by horizon \(H\) \\
``Which slot does it go into?'' & either \(x_i \in \{1, \dots, K\}\) \emph{or} \(y_{i,k} \in \{0,1\}\) --- index \emph{vs.} one-hot \\
``Which edge is used?'' & \(x_{ij} \in \{0,1\} \;\forall (i,j) \in E\) --- one Boolean per edge \\
\end{symboltable}


The fourth row deserves attention because it is a common modeling
decision: when the answer is ``which of \(K\) options,'' do you use a
single integer variable taking values in \(\{1, \dots, K\}\), or \(K\)
Boolean variables that sum to 1? Both can be correct. They differ in the
kinds of constraints they make easy to write and in the structure a
solver can exploit.

\begin{platypusinfo}
Real workforce-scheduling systems often use more
aggregate variables than the simple worker--shift assignment variable
above. In a column-generation model, for example, a variable may
represent an entire shift pattern for one worker over the planning
horizon. Since the number of possible patterns is usually enormous, the
model is solved iteratively: start with a manageable subset of patterns,
solve the restricted model, then generate additional useful patterns
through a separate pricing problem. That is a powerful technique, but it
is a different modeling level from the first paper formulation developed
here.
\end{platypusinfo}

\subsection{Step 4 --- Constraints encode validity}\label{step-4-constraints-encode-validity}

Constraints are usually the most complex part of the model. For simple
textbook examples, it may be enough to translate each ``must'' sentence
directly into one formula. But even the scenario-based knapsack example
showed that this is not always the whole story: sometimes a clear paper
statement needs auxiliary variables and additional constraints before it
becomes solver-friendly. In richer applications, one must also
distinguish between hard constraints, which define feasibility, and soft
constraints, which may be violated at a penalty. The present example is
cleanly specified, so we can ignore soft constraints for now.

We now identify the constraint families and formulate them one by one.
For hard constraints, each new rule narrows the feasible region: it
removes assignments that are not allowed, but it does not make
previously invalid assignments valid. This makes constraints relatively
modular. Still, a single rule may require several formulas, and part of
a rule may already be enforced by earlier constraints or by variable
domains. In this example, however, the mapping is pleasantly direct:
four business rules become four constraint families.

\paragraph{Coverage.} Every shift must be staffed with exactly the
headcount it requires.

\[
\sum_{w \in W} x_{w,s} \;=\; n_s \qquad \forall s \in S.
\]

This is one equation per shift. The sum on the left counts the number of
workers assigned to \(s\), and the right-hand side is the demanded
headcount. The \(\forall s \in S\) is doing real work: it makes this not
a single equation but a family of \(|S|\) equations.

\paragraph{Availability.} A worker may not be assigned to a shift they are
not available for.

\[
x_{w,s} \;\le\; a_{w,s} \qquad \forall w \in W, \; s \in S.
\]

If \(a_{w,s} = 0\), the inequality forces \(x_{w,s} = 0\). If
\(a_{w,s} = 1\), the inequality becomes \(x_{w,s} \le 1\), which is
already implied by the variable's domain, so the constraint is silent.

\paragraph{Skill match.} If a worker is assigned to a shift, the worker
must hold the skill the shift requires.

\[
x_{w,s} = 1 \;\Rightarrow\; r_s \in h_w \qquad \forall w \in W, \; s \in S.
\]

This is a clear paper-level statement, but it is not the most useful
solver-facing form because the right-hand side is a fact about the input
data, not a decision. The direct rewrite is:

\[
x_{w,s} \;\le\; [r_s \in h_w] \qquad \forall w \in W, \; s \in S,
\]

where the indicator on the right is \lstinline!1! if the
skill requirement is met and \lstinline!0! otherwise. The
compatibility values are computed from the data before the solver is
called; incompatible pairs are forced to zero, compatible pairs are left
to the other constraints.

\paragraph{No double-booking.} No worker may be assigned to more than one
shift on the same day.

\[
\sum_{s \in S_t} x_{w,s} \;\le\; 1 \qquad \forall w \in W, \; t \in T.
\]

For every worker--day combination, the sum runs only over the shifts
that fall on that day and requires that no more than one of them be
assigned to this worker. This is a typical use of a derived parameter
(\(S_t\)) to keep the constraint readable.

\subsection{Step 5 --- The objective picks among valid choices}\label{step-5-the-objective-picks-among-valid-choices}

We have specified what counts as a valid roster. Now we say which valid
roster we prefer.

\paragraph{Total cost.} Minimize the total cost of the assignments we make.

\[
\min \;\sum_{w \in W} \sum_{s \in S} c_{w,s} \, x_{w,s}.
\]

The sum runs over all worker--shift pairs, but only the pairs with
\(x_{w,s} = 1\) contribute.

Two notes on objectives. First, the sense ---
\lstinline!min! or \lstinline!max! --- is
part of the specification; ``optimize \(f\)'' is incomplete. Second, an
objective is a single expression. If the problem talks about several
goals at once --- cost \emph{and} overtime \emph{and} fairness --- you
must combine them yourself, by weighted sum or lexicographic ordering.
There is no implicit averaging. And a model can have no objective at
all: pure feasibility (``is \emph{any} valid roster possible?'') is a
legitimate question, and forcing an artificial objective onto it is a
mistake.
